\ELhold

\ELsection{سرفصل‌ها}{سرفصل‌ها}

\begin{ELunit}

\ELfigure{m03-course-outline}{}

\end{ELunit}

\ELhold

\ELsection{حل مسائل الکتریسیته ساکن}{حل مسائل الکتریسیته ساکن}

\begin{ELunit}

\begin{itemize}
\tightlist
\item
  مسائل الکتریسیته ساکن مسائلی هستند که با تأثیرات بارهای الکتریکی در حال سکون سروکار دارند. حل این نوع مسائل معمولاً شامل تعیین پتانسیل الکتریکی، شدت میدان الکتریکی و یا توزیع بار الکتریکی است.
\item
  در صورتیکه توزیع بار الکتریکی داده شده باشد، با استفاده از روابط فصل قبل هم پتانسیل و هم میدان الکتریکی قابل محاسبه خواهد بود.

  \begin{itemize}
  \tightlist
  \item
    در بسیاری از مسائل عملی توزیع دقیق بار در تمام نقاط معلوم نیست و یافتن پتانسیل و میدان الکتریکی مستقیماً با استفاده از فرمول‌های فصل قبل امکانپذیر نمی‌باشد. \ELim{\ELto} امکان استفاده از روش تصاویر در دسته ای خاص از مسائل
  \item
    در نوع دیگری از مسائل، پتانسیل تمام اجسام هادی معلوم است و می‌خواهیم پتانسیل و شدت میدان الکتریکی در فضای پیرامون و نیز توزیع بارهای سطحی را روی مرزهای هادی پیدا کنیم. \ELim{\ELto} حل مسائل مقدار مرزی
  \end{itemize}
\end{itemize}

\ELfigure{m03-map}{}

\begin{itemize}
\tightlist
\item
  فرم معادلات پواسون و لاپلاس که بر پتانسیل الکتریکی حاکم هستند چگونه است و هرکدام در چه حالتی کاربرد دارند؟
\item
  چگونه می توان با حل معادلات پواسون و لاپلاس، پتانسیل الکتریکی و شدت میدان الکتریکی را بدست آورد؟
\item
  آیا جواب معادله پواسون که شرایط مرزی مشخصی را برآورده می کند یکتاست؟
\end{itemize}

\end{ELunit}

\ELhold

\ELsection{معادلات پواسون و لاپلاس}{معادلات پواسون و لاپلاس}

\begin{ELunit}

\begin{itemize}
\tightlist
\item
  دو معادله اصلی حاکم بر الکتریسیته ساکن در کلیه محیط‌ها
\end{itemize}

\ELdisp{\nabla\cdot\vect{D}\ELbr =\rho\qquad\ELbr \qquad\ELbr  \nabla\times\vect{E}\ELbr =0}

\begin{itemize}
\tightlist
\item
  از طرفی داریم
\end{itemize}

\ELdisp{\vect{E}\ELbr =-\nabla V}
\ELdisp{\vect{D}\ELbr =\epsilon\vect{E}\quad\ELbr \ELbr \Longrightarrow\quad\ELbr \nabla\cdot\epsilon\vect{E}\ELbr =\rho\quad\ELbr \ELbr \Longrightarrow\quad\ELbr \nabla\cdot(\epsilon\nabla V)\ELbr =-\rho}

\end{ELunit}

\begin{ELjoin}

\begin{itemize}
\tightlist
\item
  اگر محیط همگن باشد (\ELim{\epsilon} ثابت باشد)
\end{itemize}

\begin{ELimportant}{}

\ELdisp{\nabla^2V\ELbr =-\frac{\rho}{\epsilon}}

\end{ELimportant}

\end{ELjoin}

\begin{ELunit}

\begin{itemize}
\item
  \ELim{\rho}: چگالی بار آزاد
\item
  در معادله فوق که به \textbf{معادله پواسون} معروف است، از عملگر \textbf{لاپلاسین} (دیورژانسِ گرادیان) استفاده شده است.
\item
  لاپلاسین یک کمیت عددی در دستگاه‌های مختصات مختلف

  \begin{itemize}
  \tightlist
  \item
    کارتزین
    \ELdisp{\nabla^2V\ELbr =\frac{\partial^2V}{\partial x^2}+\frac{\partial^2V}{\partial y^2}+\frac{\partial^2V}{\partial z^2}}
  \item
    استوانه ای
    \ELdisp{\nabla^2V\ELbr =\frac1r\frac{\partial}{\partial r}\left(r\frac{\partial V}{\partial r}\right)+\frac{1}{r^2}\frac{\partial^2V}{\partial\phi^2}+\frac{\partial^2V}{\partial z^2}}
  \item
    کروی
    \ELdisp{\nabla^2V\ELbr =\frac{1}{R^2}\frac{\partial}{\partial R}\left(R^2\frac{\partial V}{\partial R}\right)+\frac{1}{R^2\sin\theta}\frac{\partial}{\partial\theta}\left(\sin\theta\frac{\partial V}{\partial\theta}\right)+\frac{1}{R^2\sin^2\theta}\frac{\partial^2V}{\partial\phi^2}}
  \end{itemize}
\item
  حل معادله پواسون در فضای سه بعدی و با شرایط مرزی مشخص معمولاً کار ساده‌ای نیست.
\item
  در نقاطی از یک محیط ساده که بار آزادی وجود ندارد
\end{itemize}

\end{ELunit}

\begin{ELimportant}{}

\ELdisp{\nabla^2V\ELbr =0}

\end{ELimportant}

\begin{ELunit}

\begin{itemize}
\tightlist
\item
  معادله فوق \textbf{معادله لاپلاس} نامیده می‌شود.

  \begin{itemize}
  \tightlist
  \item
    این معادله حاکم بر مسائلی است که شامل مجموعه‌ای از هادی‌ها باشد که در پتانسیل‌های مختلف نگه داشته شوند.
  \item
    وقتی با استفاده از این معادله \ELim{V} بدست آمد، \ELim{\vect{E}} و توزیع بار روی سطوح هادی نیز به سادگی بدست خواهند آمد
  \end{itemize}
\end{itemize}

\ELdisp{\vect{E}\ELbr =-\nabla V\qquad\ELbr \qquad\ELbr  \rho_s\ELbr =-\epsilon E_n}

\end{ELunit}

\ELnextnum{3-1}

\begin{ELexample}{}

دو صفحه خازن صفحه ای موازی، به فاصله \ELim{d} از یکدیگر قرار دارند و مطابق شکل زیر در پتانسیل های \ELim{0} و \ELim{V_0} نگه داشته می شوند. با چشم پوشی از میدان های حاشیه ای مطلوبست

\begin{itemize}
\tightlist
\item
  الف. پتانسیل در تمام نقاط بین صفحات
\item
  ب. چگالی های بار سطحی روی صفحات
\end{itemize}

\ELfigure{m03-ex1}{}

\begin{ELsolution}

\begin{itemize}
\tightlist
\item
  الف. استفاده از معادله لاپلاس
\end{itemize}

\ELdisp{\nabla^2V\ELbr =0\quad\ELbr \ELbr \Longrightarrow\quad\ELbr \frac{d^2V}{dy^2}\ELbr =0\quad\ELbr \ELbr \Longrightarrow\quad\ELbr \frac{dV}{dy}\ELbr =C_1\quad\ELbr \ELbr \Longrightarrow\quad\ELbr  V\ELbr =C_1y+C_2}

\begin{itemize}
\tightlist
\item
  شرایط مرزی:
\end{itemize}

\ELdisp{\left.\begin{aligned}&\text{\lr{At }}y=0,\quad V=0\\&\text{\lr{At }}y=d,\quad V=V_0\end{aligned}\right\}\quad\ELbr \ELbr \Longrightarrow\quad\ELbr  V\ELbr =\frac{V_0}{d}y}

\begin{itemize}
\tightlist
\item
  ب.
\end{itemize}

\ELdisp{\vect{E}\ELbr =-\nabla V\quad\ELbr \ELbr \Longrightarrow\quad\ELbr \vect{E}\ELbr =-\uvec{y}\frac{dV}{dy}\ELbr =-\uvec{y}\frac{V_0}{d}}
\ELdisp{E_n\ELbr =\uvec{n}\cdot\vect{E}\ELbr =\frac{\rho_s}{\epsilon}}

\begin{itemize}
\tightlist
\item
  برای صفحه پایینی
\end{itemize}

\ELdisp{\uvec{n}\ELbr =\uvec{y},\qquad\ELbr  E_{n\ell}\ELbr =-\frac{V_0}{d},\qquad\ELbr  \rho_{s\ell}\ELbr =-\frac{\epsilon V_0}{d}}

\begin{itemize}
\tightlist
\item
  برای صفحه بالایی
\end{itemize}

\ELdisp{\uvec{n}\ELbr =-\uvec{y},\qquad\ELbr  E_{nu}\ELbr =\frac{V_0}{d},\qquad\ELbr  \rho_{su}\ELbr =\frac{\epsilon V_0}{d}}

\end{ELsolution}

\end{ELexample}

\ELnextnum{3-2}

\begin{ELexample}{}

با حل معادله پواسون و لاپلاس نسبت به \ELim{V}، میدان \ELim{\vect{E}} را درون و بیرون یک ابر الکترونی کروی با چگالی بار حجمی یکنواخت \ELim{\rho=-\rho_0} در \ELim{0\leq R\leq b} و \ELim{\rho=0} در \ELim{R>b} تعیین کنید.

\begin{ELsolution}

\begin{itemize}
\tightlist
\item
  درون ابر الکترونی \ELim{\ELto} استفاده از معادله پواسون
\end{itemize}

\ELdisp{\nabla^2V\ELbr =-\frac\rho\epsilon\quad\ELbr \ELbr \Longrightarrow\quad\ELbr \frac{1}{R^2}\frac{d}{dR}\left(R^2\frac{dV_i}{dR}\right)\ELbr =\frac{\rho_0}{\epsilon_0}\quad\ELbr \ELbr \Longrightarrow\quad\ELbr \frac{d}{dR}\left(R^2\frac{dV_i}{dR}\right)\ELbr =\frac{\rho_0}{\epsilon_0}R^2\quad\ELbr \ELbr \Longrightarrow\quad\ELbr  R^2\frac{dV_i}{dR}\ELbr =\frac{\rho_0}{3\epsilon_0}R^3+C_1}
\ELdisp{\frac{dV_i}{dR}\ELbr =\frac{\rho_0}{3\epsilon_0}R+\frac{C_1}{R^2}\quad\ELbr \ELbr \Longrightarrow\quad\ELbr \vect{E}_i\ELbr =-\nabla V_i\ELbr =-\uvec{R}\left(\frac{dV_i}{dR}\right)}

\begin{itemize}
\tightlist
\item
  با توجه به اینکه میدان در \ELim{R=0} نمی تواند بی نهایت باشد
\end{itemize}

\ELdisp{C_1\ELbr =0\quad\ELbr \ELbr \Longrightarrow\quad\ELbr \vect{E}_i\ELbr =-\uvec{R}\frac{\rho_0}{3\epsilon_0}R,\qquad\ELbr  0\leq R\leq b}

\begin{itemize}
\tightlist
\item
  بیرون ابر الکترونی \ELim{\ELto} استفاده از معادله لاپلاس
\end{itemize}

\ELdisp{\nabla^2V\ELbr =0\quad\ELbr \ELbr \Longrightarrow\quad\ELbr \frac{1}{R^2}\frac{\partial}{\partial R}\left(R^2\frac{dV_o}{dR}\right)\ELbr =0\quad\ELbr \ELbr \Longrightarrow\quad\ELbr \frac{dV_o}{dR}\ELbr =\frac{C_2}{R^2}}
\ELdisp{\vect{E}_o\ELbr =-\nabla V_o\ELbr =-\uvec{R}\frac{dV_o}{dR}\ELbr =-\uvec{R}\frac{C_2}{R^2}}

\begin{itemize}
\tightlist
\item
  با توجه به پیوستگی محیط، میدان های \ELim{\vect{E}_i} و \ELim{\vect{E}_o} در \ELim{R=b} با هم برابرند
\end{itemize}

\ELdisp{\frac{C_2}{b^2}\ELbr =\frac{\rho_0}{3\epsilon_0}b\quad\ELbr \ELbr \Longrightarrow\quad\ELbr  C_2\ELbr =\frac{\rho_0b^3}{3\epsilon_0}\quad\ELbr \ELbr \Longrightarrow\quad\ELbr \vect{E}_o\ELbr =-\uvec{R}\frac{\rho_0b^3}{3\epsilon_0R^2},\qquad\ELbr  R\geq b}

\end{ELsolution}

\end{ELexample}

\ELhold

\ELsection{قضیه یکتایی}{قضیه یکتایی}

\begin{ELunit}

\ELfigure{m03-map-unique}{}

\begin{itemize}
\tightlist
\item
  یک جواب معادله پواسون (و در حالت خاص معادله لاپلاس) که شرایط مرزی مفروضی را برآورده می‌کند، جواب یکتای معادله مذکور خواهد بود.
\item
  نتیجه این قضیه آن است که یک جواب هر مسئله الکتریسیته ساکن که در شرایط مرزی صدق کند، بدون توجه به روش بدست آوردن آن جواب، تنها جواب ممکن است.
\end{itemize}

\end{ELunit}

\ELhold

\ELsection{روش تصاویر}{روش تصاویر}

\begin{ELunit}

\ELfigure{m03-map-images}{}

\begin{itemize}
\item
  روش تصاویر چیست و چگونه می توان از این روش برای حل دسته ای خاص از مسائل الکتریسیته ساکن استفاده کرد؟
\item
  در دسته‌ای از مسائل الکترومغناطیس، حل معادله لاپلاس به همراه برقراری شرایط مرزی بسیار مشکل است. اما در این مسائل می‌توان شرایط مرزی را به کمک \textbf{بارهای تصویر} ارضا کرد.

  \begin{itemize}
  \tightlist
  \item
    به این روش که عبارتست از جایگزینی سطوح مرزی با بارهای تصویرِ معادل \textbf{روش تصویر} گویند.
  \end{itemize}
\item
  به عنوان مثال مسئله زیر را در نظر می‌گیریم. می‌خواهیم پتانسیل را در تمام نقط \ELim{y>0} بدست آوریم.
\end{itemize}

\ELfigure{m03-image-problem}{}

\begin{itemize}
\tightlist
\item
  از قانون گوس نمی‌توان برای محاسبه میدان و سپس پتانسیل استفاده کرد زیرا تعیین سطح گوسی ممکن نیست.
\item
  به صورت مستقیم نیز نمی توان پتانسیل را محاسبه کرد زیرا حضور بار مثبت \ELim{Q} باعث القای بارهای منفی روی سطح هادی می‌شود و چگالی بار سطحی \ELim{\rho_s} را نتیجه می‌دهد. داریم
\end{itemize}

\ELdisp{V(x,y,z)\ELbr =\frac{Q}{4\pi\epsilon_0\sqrt{x^2+(y-d)^2+z^2}}+\frac{1}{4\pi\epsilon_0}\int_S\frac{\rho_s}{R_1}\,ds}

\begin{itemize}
\tightlist
\item
  مشکل اینجاست که ابتدا باید \ELim{\rho_s} را بدست آوریم و حتی در صورت بدست آوردن \ELim{\rho_s} محاسبه انتگرال سطحی مشکل خواهد بود.
\item
  حل معادله لاپلاس برای دستیابی به جواب نیز غیرممکن یا بسیار مشکل است. در واقع جواب معادله لاپلاس باید شرایط زیر را برآورده کند
\end{itemize}

\ELdisp{V(x,0,z)\ELbr =0}
\ELdisp{V\to\frac{Q}{4\pi\epsilon_0R},\ \text{\lr{as}}\ R\to0}

\begin{itemize}
\tightlist
\item
  در فواصل بسیار دور از بار \ELim{Q}، پتانسیل باید صفر باشد
\end{itemize}

\ELdisp{V(x,y,z)\ELbr =V(-x,y,z)}
\ELdisp{V(x,y,z)\ELbr =V(x,y,-z)}

\begin{itemize}
\tightlist
\item
  می توان چنین مسئله ای را به سادگی با استفاده از روش تصاویر حل کرد.
\end{itemize}

\end{ELunit}

\ELhold

\ELsubsection{بار نقطه ای و صفحه هادی}{بار نقطه ای و صفحه هادی}

\begin{ELunit}

\begin{itemize}
\tightlist
\item
  صفحه هادی را برمی‌داریم و آن را با بار نقطه‌ای تصویر \ELim{-Q} در \ELim{y=-d} جایگزین می‌کنیم
\end{itemize}

\ELfigure{m03-image-plane}{}

\end{ELunit}

\begin{ELimportant}{}

\ELdisp{V(x,y,z)\ELbr =\frac{Q}{4\pi\epsilon_0}\left(\frac{1}{R_+}-\frac{1}{R_-}\right)}
\ELdisp{R_+\ELbr =\left[x^2+(y-d)^2+z^2\right]^{1/2},\qquad\ELbr  R_-\ELbr =\left[x^2+(y+d)^2+z^2\right]^{1/2}}

\end{ELimportant}

\begin{ELunit}

\begin{itemize}
\tightlist
\item
  به سادگی می‌توان نشان داد که این جواب در معادله لاپلاس صدق کرده و شرایط مرزی را ارضا می‌کند. بنابراین این جواب یک جواب مسئله بوده و مطابق با قضیه یکتایی تنها جواب آن نیز می‌باشد.
\item
  با بدست آوردن پتانسیل، محاسبه میدان و توزیع بار سطحی روی هادی آسان خواهد بود.
\end{itemize}

\end{ELunit}

\begin{ELremark}{}

نکته‌ای که باید توجه شود این است که این روش نمی‌تواند برای محاسبه پتانسیل در نواحی \ELim{y<0} بکار رود.

\end{ELremark}

\ELhold

\ELsubsection{بار خطی و استوانه هادی موازی}{بار خطی و استوانه هادی موازی}

\begin{ELunit}

\ELfigure{m03-line-cylinder}{}

\begin{itemize}
\tightlist
\item
  هدف یافتن پتانسیل در خارج استوانه هادی است.
\item
  تصویر باید یک بار خطی موازی در درون استوانه باشد تا سطح استوانه در \ELim{r=a} را یک سطح هم‌پتانسیل کند. هدف یافتن \ELim{d_i} و \ELim{\rho_i} است.
\item
  فرض می‌کنیم
\end{itemize}

\end{ELunit}

\begin{ELimportant}{}

\ELdisp{\rho_i\ELbr =-\rho_\ell}

\end{ELimportant}

\begin{ELunit}

\begin{itemize}
\tightlist
\item
  پتانسیل در فاصله \ELim{r} ناشی از بار خطی به چگالی \ELim{\rho_\ell} با فرض مرجع در \ELim{r_0} برابر است با
\end{itemize}

\ELdisp{V\ELbr =-\int_{r_0}^{r}E_r\,dr\ELbr =-\frac{\rho_\ell}{2\pi\epsilon_0}\int_{r_0}^{r}\frac1r\,dr\ELbr =\frac{\rho_\ell}{2\pi\epsilon_0}\ln\frac{r_0}{r}}

\begin{itemize}
\tightlist
\item
  پتانسیل در نقطه \ELim{M} روی سطح استوانه برابر است با
\end{itemize}

\ELdisp{\begin{aligned}V_M&=\frac{\rho_\ell}{2\pi\epsilon_0}\ln\frac{r_0}{r}-\frac{\rho_\ell}{2\pi\epsilon_0}\ln\frac{r_0}{r_i}\\&=\frac{\rho_\ell}{2\pi\epsilon_0}\ln\frac{r_i}{r}\end{aligned}}

\begin{itemize}
\tightlist
\item
  بنابراین سطوح هم‌پتانسیل به صورت زیر مشخص می‌شوند
\end{itemize}

\ELdisp{\frac{r_i}{r}\ELbr =\text{\lr{Constant}}}

\ELfigure{m03-similar-triangles}{}

\end{ELunit}

\begin{ELjoin}

\begin{itemize}
\tightlist
\item
  اگر شرط \ELim{\dfrac{d_i}{a}=\dfrac ad} برقرار باشد، از تشابه بین دو مثلث \ELim{OMP_i} و \ELim{OPM} همواره نسبت \ELim{\dfrac{r_i}{r}} ثابت خواهد بود. بنابراین
\end{itemize}

\begin{ELimportant}{}

\ELdisp{d_i\ELbr =\frac{a^2}{d}}

\end{ELimportant}

\end{ELjoin}

\begin{ELjoin}

\begin{itemize}
\tightlist
\item
  به این ترتیب بار خطی \ELim{-\rho_\ell} می‌تواند جایگزین سطح استوانه‌ای هادی شود و پتانسیل و میدان الکتریکی را هر نقطه خارج استوانه بدست آورد.
\end{itemize}

\begin{ELremark}{}

نکته‌ای که باید توجه شود این است که این روش نمی‌تواند برای محاسبه پتانسیل در نواحی درون استوانه هادی بکار رود.

\end{ELremark}

\end{ELjoin}

\ELnextnum{3-3}

\begin{ELexample}{}

ظرفیت در واحد طول بین دو سیم هادی دایروی موازی بسیار بلند به شعاع \ELim{a} و فاصله \ELim{D} را بدست آورید.

\ELfigure{m03-ex3}{}

\begin{ELsolution}

\begin{itemize}
\tightlist
\item
  بارهای \ELim{+Q} و \ELim{-Q} را به ترتیب روی هادی های ۱ و ۲ قرار می دهیم.
\item
  با استفاده از روش تصاویر می توان سطوح هم پتانسیل سیم ها را با دو بارخطی \ELim{+\rho_\ell} و \ELim{-\rho_\ell} جایگزین کرد.
\end{itemize}

\ELdisp{V_2\ELbr =\frac{\rho_\ell}{2\pi\epsilon_0}\ln\frac ad}
\ELdisp{V_1\ELbr =-\frac{\rho_\ell}{2\pi\epsilon_0}\ln\frac ad}
\ELdisp{C\ELbr =\frac{\rho_\ell}{V_1-V_2}\ELbr =\frac{\pi\epsilon_0}{\ln(d/a)}}
\ELdisp{d\ELbr =D-d_i\ELbr =D-\frac{a^2}{d}\quad\ELbr \ELbr \Longrightarrow\quad\ELbr  d\ELbr =\tfrac12\left(D+\sqrt{D^2-4a^2}\right)}
\ELdisp{C\ELbr =\frac{\pi\epsilon_0}{\ln\left[(D/2a)+\sqrt{(D/2a)^2-1}\right]}\qquad\ELbr (\mathrm{F/m})}

\begin{ELimportant}{}

\ELdisp{C\ELbr =\frac{\pi\epsilon_0}{\cosh^{-1}(D/2a)}\qquad\ELbr (\mathrm{F/m})}

\end{ELimportant}

\end{ELsolution}

\end{ELexample}

\ELhold

\ELsection{مسائل مقدار مرزی}{مسائل مقدار مرزی}

\begin{ELunit}

\ELfigure{m03-map-bvp}{}

\begin{itemize}
\tightlist
\item
  روش تصاویر در حل مسائلی که شامل بار آزاد در نزدیکی مرزهای هادی با شکل هندسی ساده هستند، کاربرد دارد.
\item
  اگر مسئله شامل دسته‌ای از هادی‌های نگه داشته شده در پتانسیل مشخص و بدون بار آزاد مجزا باشد، با روش تصاویر قابل حل نخواهد بود و نیاز به حل معادله لاپلاس است.
\item
  معادله لاپلاس، یک معادله دیفرانسیل جزئی است. مسائلی را که از معادلات دیفرانسیل جزئی با شرایط مرزی مشخص پیروی می‌کنند، \textbf{مسائل مقدار مرزی} گویند.
\item
  مسائل مقدار مرزی در توابع پتانسیل به دسته تقسیم می‌شوند

  \begin{itemize}
  \tightlist
  \item
    مسائل دیریشله: مقدار پتانسیل در تمام نقاط مرزها مشخص شده است.
  \item
    مسائل نویمن: مشتق عمودی پتانسیل در تمام مرزها مشخص است.
  \item
    مسائل مقدار مرزی ترکیبی: پتانسل در بعضی از نقاط مرزها و مشتق عمودی پتانسیل در بقیه نقاط مشخص شده است.
  \end{itemize}
\item
  نحوه حل مسائل مقدار مرزی در حوزه الکتریسیته ساکن در دستگاه های مختصات مختلف و برای دستیابی به توزیع پتانسیل چگونه است؟
\end{itemize}

\end{ELunit}

\ELhold

\ELsubsection{مختصات کارتزین}{مختصات کارتزین}

\begin{ELunit}

\begin{itemize}
\tightlist
\item
  معادله لاپلاس در مختصات کارتزین
\end{itemize}

\ELdisp{\frac{\partial^2V}{\partial x^2}+\frac{\partial^2V}{\partial y^2}+\frac{\partial^2V}{\partial z^2}\ELbr =0}

\begin{itemize}
\tightlist
\item
  با استفاده از روش جداسازی متغیرها
\end{itemize}

\ELdisp{V(x,y,z)\ELbr =X(x)Y(y)Z(z)}
\ELdisp{Y(y)Z(z)\frac{d^2X(x)}{dx^2}+X(x)Z(z)\frac{d^2Y(y)}{dy^2}+X(x)Y(y)\frac{d^2Z(z)}{dz^2}\ELbr =0}
\ELdisp{\underbrace{\frac{1}{X(x)}\frac{d^2X(x)}{dx^2}}_{-k_x^2}+\underbrace{\frac{1}{Y(y)}\frac{d^2Y(y)}{dy^2}}_{-k_y^2}+\underbrace{\frac{1}{Z(z)}\frac{d^2Z(z)}{dz^2}}_{-k_z^2}\ELbr =0}
\ELdisp{k_x^2+k_y^2+k_z^2\ELbr =0}

\end{ELunit}

\begin{ELimportant}{}

\ELdisp{\frac{d^2X(x)}{dx^2}+k_x^2X(x)\ELbr =0\qquad\ELbr  \frac{d^2Y(y)}{dy^2}+k_y^2Y(y)\ELbr =0\qquad\ELbr  \frac{d^2Z(z)}{dz^2}+k_z^2Z(z)\ELbr =0}

\end{ELimportant}

\textbf{\lr{Possible Solutions of $X''(x)+k_x^2X(x)=0$}}

\Needspace{10\baselineskip}
\begin{ELltr}
\begin{longtable}[c]{llll}
\toprule\noalign{}
\ELtablehead \ELim{k_x^2} & \ELim{k_x} & \ELim{X(x)} & \lr{Exponential forms of} \ELim{X(x)} \\\noalign{\vskip 4pt}
\midrule\noalign{}
\endhead
\bottomrule\noalign{}
\endlastfoot
\ELim{0} & \ELim{0} & \ELim{A_0x+B_0} & \\\noalign{\vskip 4pt}
\ELim{+} & \ELim{k} & \ELim{A_1\sin kx+B_1\cos kx} & \ELim{C_1e^{jkx}+D_1e^{-jkx}} \\\noalign{\vskip 4pt}
\ELim{-} & \ELim{jk} & \ELim{A_2\sinh kx+B_2\cosh kx} & \ELim{C_2e^{kx}+D_2e^{-kx}} \\\noalign{\vskip 4pt}
\end{longtable}\end{ELltr}


\ELnextnum{3-4}

\begin{ELexample}{}

دو صفحه هادی موازی زمین شده نیمه بینهایت به فاصله \ELim{d} از هم قرار دارند. صفحه سوم عمود بر این دو صفحه و جدا از آن ها در پتانسیل \ELim{V_0} نگه داشته می شود. توزیع پتانسیل را در ناحیه احاطه شده توسط این صفحات تعیین کنید.

\ELfigure{m03-ex4}{}

\begin{ELsolution}

\begin{itemize}
\tightlist
\item
  شرایط مرزی
\end{itemize}

\ELdisp{\text{\lr{(1)}}\ \ V(x,y,z)\ELbr =V(x,y)\qquad\ELbr \text{\lr{(2)}}\ \ V(0,y)\ELbr =V_0\qquad\ELbr \text{\lr{(3)}}\ \ V(\infty,y)\ELbr =0\qquad\ELbr \text{\lr{(4)}}\ \ V(x,0)\ELbr =0\qquad\ELbr \text{\lr{(5)}}\ \ V(x,b)\ELbr =0}

\ELdisp{\text{\lr{(1)}}\quad\ELbr \ELbr \Longrightarrow\quad\ELbr \frac{\partial V}{\partial z}\ELbr =0\quad\ELbr \ELbr \Longrightarrow\quad\ELbr  Z(z)\ELbr =B_0}
\ELdisp{\left.\begin{aligned}k_z^2&=-\frac{1}{Z(z)}\frac{d^2Z(z)}{dz^2}=0\\k_x^2+k_y^2+k_z^2&=0\end{aligned}\right\}\quad\ELbr  k_y^2\ELbr =-k_x^2\ELbr =k^2\qquad\ELbr (k\ \text{\lr{real}})}
\ELdisp{k_x\ELbr =jk\quad\ELbr \ELbr \Longrightarrow\quad\ELbr  X(x)\ELbr =C_2e^{kx}+D_2e^{-kx}\quad\ELbr \overset{(3)}{\Longrightarrow}\quad\ELbr  X(x)\ELbr =D_2e^{-kx}}
\ELdisp{k_y\ELbr =k\quad\ELbr \ELbr \Longrightarrow\quad\ELbr  Y(y)\ELbr =A_1\sin ky+B_1\cos ky\quad\ELbr \overset{(4)}{\Longrightarrow}\quad\ELbr  Y(y)\ELbr =A_1\sin ky}
\ELdisp{\begin{aligned}V_n(x,y)&=(B_0D_2A_1)e^{-kx}\sin ky\\&=C_ne^{-kx}\sin ky\end{aligned}}
\ELdisp{\overset{(5)}{\Longrightarrow}\quad\ELbr  V_n(x,b)\ELbr =C_ne^{-kx}\sin kb\ELbr =0\quad\ELbr \ELbr \Longrightarrow\quad\ELbr \sin kb\ELbr =0\quad\ELbr \ELbr \Longrightarrow\quad\ELbr  kb\ELbr =n\pi}
\ELdisp{k\ELbr =\frac{n\pi}{b},\qquad\ELbr  n\ELbr =1,2,3,\ldots}
\ELdisp{V_n(x,y)\ELbr =C_ne^{-n\pi x/b}\sin\frac{n\pi}{b}y}

\begin{itemize}
\tightlist
\item
  رابطه فوق در معادله لاپلاس صدق می کند اما به تنهایی نمی تواند شرط مرزی دوم در \ELim{x=0} را برای همه \ELim{y} ها از \ELim{0} تا \ELim{b} ارضا کند.
\item
  از آنجا که معادله لاپلاس یک معادله دیفرانسیل جزیی خطی است، ترکیب \ELim{V_n} هایی به فرم فوق و به ازای مقادیر مختلف \ELim{n} نیز یک پاسخ خواهد بود.
\end{itemize}

\ELdisp{V(x,y)\ELbr =\sum_{n=1}^{\infty}V_n(x,y)}

\ELdisp{\overset{(2)}{\Longrightarrow}\quad\ELbr \begin{aligned}V(0,y)&=\sum_{n=1}^{\infty}V_n(0,y)=\sum_{n=1}^{\infty}C_n\sin\frac{n\pi}{b}y\\&=V_0,\qquad 0<y<b\end{aligned}}

\begin{itemize}
\tightlist
\item
  رابطه فوق بسط سری فوریه یک تابع متناوب و فرد است که در بازه \ELim{0<y<b} مقدار آن \ELim{V_0} می باشد.
\end{itemize}

\ELfigure{m03-ex4-square}{}

\begin{ELremark}{یادآوری}

\ELdisp{f(t)\ELbr =a_0+\sum_{n=1}^{\infty}\left[a_n\cos\left(\frac{2n\pi}{T}t\right)+b_n\sin\left(\frac{2n\pi}{T}t\right)\right]}
\ELdisp{b_n\ELbr =\frac2T\int_{-T/2}^{T/2}f(t)\sin\left(\frac{2n\pi}{T}t\right)dt\qquad\ELbr  a_n\ELbr =\frac2T\int_{-T/2}^{T/2}f(t)\cos\left(\frac{2n\pi}{T}t\right)dt\qquad\ELbr  a_0\ELbr =\frac1T\int_{-T/2}^{T/2}f(t)\,dt}

\end{ELremark}

\ELdisp{\begin{aligned}C_n&=\frac{2}{2b}\left[\int_0^bV_0\sin\left(\frac{2n\pi}{2b}y\right)dy+\int_b^{2b}-V_0\sin\left(\frac{2n\pi}{2b}y\right)dy\right]\\&=\frac{V_0}{n\pi}(-2\cos n\pi+2)=\begin{cases}\dfrac{4V_0}{n\pi}&\text{\lr{if }}n\text{\lr{ is odd}}\\\noalign{\vskip 2mm}0&\text{\lr{if }}n\text{\lr{ is even}}\end{cases}\end{aligned}}

\begin{ELimportant}{}

\ELdisp{\begin{aligned}V(x,y)&=\sum_{n=1}^{\infty}C_ne^{-n\pi x/b}\sin\frac{n\pi}{b}y\\&=\frac{4V_0}{\pi}\sum_{n=\text{\lr{odd}}}^{\infty}\frac1ne^{-n\pi x/b}\sin\frac{n\pi}{b}y,\end{aligned}}
\ELdisp{n\ELbr =1,3,5,\ldots,\qquad\ELbr  x>0\quad\ELbr \text{\lr{and}}\quad\ELbr  0<y<b}

\end{ELimportant}

\end{ELsolution}

\end{ELexample}

\ELhold

\ELsubsection{مختصات استوانه ای}{مختصات استوانه ای}

\begin{ELunit}

\begin{itemize}
\tightlist
\item
  معادله لاپلاس در مختصات استوانه ای
\end{itemize}

\ELdisp{\frac1r\frac{\partial}{\partial r}\left(r\frac{\partial V}{\partial r}\right)+\frac{1}{r^2}\frac{\partial^2V}{\partial\phi^2}+\frac{\partial^2V}{\partial z^2}\ELbr =0}

\begin{itemize}
\tightlist
\item
  با فرض بزرگ بودن بعد طولی در مقایسه با شعاع، در هندسه استوانه ای، پتانسیل مستقل از \ELim{z} است.
\end{itemize}

\ELdisp{\frac1r\frac{\partial}{\partial r}\left(r\frac{\partial V}{\partial r}\right)+\frac{1}{r^2}\frac{\partial^2V}{\partial\phi^2}\ELbr =0}

\begin{itemize}
\tightlist
\item
  با استفاده از روش جداسازی متغیرها
\end{itemize}

\ELdisp{V(r,\phi)\ELbr =R(r)\Phi(\phi)}
\ELdisp{\underbrace{\frac{r}{R(r)}\frac{d}{dr}\left[r\frac{dR(r)}{dr}\right]}_{k^2}+\underbrace{\frac{1}{\Phi(\phi)}\frac{d^2\Phi(\phi)}{d\phi^2}}_{-k^2}\ELbr =0}

\end{ELunit}

\begin{ELjoin}

\begin{itemize}
\tightlist
\item
  معادله دیفرانسیل برحسب \ELim{\phi}
\end{itemize}

\begin{ELimportant}{}

\ELdisp{\frac{d^2\Phi(\phi)}{d\phi^2}+k^2\Phi(\phi)\ELbr =0}

\end{ELimportant}

\end{ELjoin}

\begin{ELunit}

\begin{itemize}
\tightlist
\item
  \textbf{وقتی \ELim{k\neq0}}

  \begin{itemize}
  \tightlist
  \item
    در شکل‌بندی‌های استوانه‌ای مدور، توابع پتانسیل و در نتیجه \ELim{\Phi(\phi)} نسبت به \ELim{\phi} متناوب بوده و توابع هذلولوی بکار نمی‌روند.

    \begin{itemize}
    \tightlist
    \item
      اگر دامنه \ELim{\phi} محدود نشده باشد، \ELim{k} باید عدد صحیح باشد (\ELim{k=n})
    \end{itemize}
  \end{itemize}
\end{itemize}

\ELdisp{\Phi(\phi)\ELbr =A_\phi\sin n\phi+B_\phi\cos n\phi}

\begin{itemize}
\tightlist
\item
  معادله دیفرانسیل برحسب \ELim{r} (معادله اویلر)
\end{itemize}

\ELdisp{r^2\frac{d^2R(r)}{dr^2}+r\frac{dR(r)}{dr}-n^2R(r)\ELbr =0}
\ELdisp{R(r)\ELbr =A_rr^n+B_rr^{-n}}

\end{ELunit}

\begin{ELimportant}{}

\ELdisp{V_n(r,\phi)\ELbr =r^n(A_n\sin n\phi+B_n\cos n\phi)+r^{-n}(A'_n\sin n\phi+B'_n\cos n\phi),\qquad\ELbr  n\neq0}

\end{ELimportant}

\begin{ELunit}

\begin{itemize}
\item
  هنگامیکه ناحیه مورد نظر شامل \ELim{r=0} باشد، جملات شامل \ELim{r^{-n}} و هنگامیکه ناحیه مورد نظر شامل بی‌نهایت باشد، جملات شامل \ELim{r^n} نمی‌توانند وجود داشته باشند.
\item
  \textbf{وقتی \ELim{k=0}}
\end{itemize}

\ELdisp{\frac{d^2\Phi(\phi)}{d\phi^2}\ELbr =0}
\ELdisp{\Phi(\phi)\ELbr =A_0\phi+B_0\quad\ELbr \ELbr \Longrightarrow\quad\ELbr \Phi(\phi)\ELbr =B_0}

\begin{itemize}
\tightlist
\item
  اگر هیچ تغییرات پیرامونی وجود نداشته باشد
\end{itemize}

\ELdisp{\frac{d}{dr}\left[r\frac{dR(r)}{dr}\right]\ELbr =0}
\ELdisp{R(r)\ELbr =C_0\ln r+D_0}

\end{ELunit}

\ELnextnum{3-5}

\begin{ELexample}{}

یک کابل هم محور بسیار بلند را در نظر بگیرید. شعاع هادی داخلی \ELim{a} و در پتانسیل \ELim{V_0} نگه داشته شده است. هادی خارجی دارای شعاع \ELim{b} بوده و زمین شده است. توزیع پتانسیل در فضای بین دو هادی را بدست آورید.

\ELfigure{m03-ex5}{}

\begin{ELsolution}

\begin{itemize}
\tightlist
\item
  با توجه به بلند بودن کابل، وابستگی به \ELim{z} وجود ندارد.
\item
  با توجه به تقارن، وابستگی به \ELim{\phi} وجود ندارد (\ELim{k=0}).
\end{itemize}

\ELdisp{V(r)\ELbr =C_1\ln r+C_2}
\ELdisp{V(b)\ELbr =0,\qquad\ELbr  V(a)\ELbr =V_0}
\ELdisp{C_1\ELbr =-\frac{V_0}{\ln(b/a)},\qquad\ELbr  C_2\ELbr =\frac{V_0\ln b}{\ln(b/a)}\qquad\ELbr \ELbr \Longrightarrow\qquad\ELbr  V(r)\ELbr =\frac{V_0}{\ln(b/a)}\ln\left(\frac br\right)}

\end{ELsolution}

\end{ELexample}

\ELnextnum{3-6}

\begin{ELexample}{}

دو صفحه هادی نیمه بینهایت مطابق شکل زیر در پتانسیل های \ELim{V_0} و صفر نگه داشته شده اند. با فرض اینکه طول این صفحات در امتداد محور \ELim{z} بسیار بلند باشد، توزیع پتانسیل را در کل فضا بدست آورید.

\ELfigure{m03-ex6}{}

\begin{ELsolution}

\begin{itemize}
\tightlist
\item
  با توجه به بلند بودن صفحات، وابستگی به \ELim{z} وجود ندارد.
\item
  وابستگی به \ELim{r} وجود ندارد (\ELim{k=0}).
\end{itemize}

\ELdisp{\Phi(\phi)\ELbr =A_0\phi+B_0}
\ELdisp{V(\phi)\ELbr =\begin{cases}0&\phi=0,2\pi\\V_0&\phi=\alpha\end{cases}}

\begin{itemize}
\tightlist
\item
  \ELim{0\leq\phi\leq\alpha}
\end{itemize}

\ELdisp{V(\phi)\ELbr =\frac{V_0}{\alpha}\phi}

\begin{itemize}
\tightlist
\item
  \ELim{\alpha\leq\phi\leq2\pi}
\end{itemize}

\ELdisp{V(\phi)\ELbr =\frac{V_0}{2\pi-\alpha}(2\pi-\phi)}

\end{ELsolution}

\end{ELexample}

\ELnextnum{3-7}

\begin{ELexample}{}

یک لوله هادی استوانه ای نازک بسیار بلند به شعاع \ELim{b} به دو بخش تقسیم شده است. نیمه بالایی در پتانسیل \ELim{V_0} و نیمه پایینی در پتانسیل \ELim{-V_0} نگه داشته شده است. توزیع پتانسیل درون و بیرون لوله را بدست آورید.

\ELfigure{m03-ex7}{}

\begin{ELsolution}

\begin{itemize}
\tightlist
\item
  با توجه به بلند بودن لوله وابستگی به \ELim{z} وجود ندارد.
\end{itemize}

\ELdisp{V_n(r,\phi)\ELbr =r^n(A_n\sin n\phi+B_n\cos n\phi)+r^{-n}(A'_n\sin n\phi+B'_n\cos n\phi)}
\ELdisp{V(b,\phi)\ELbr =\begin{cases}V_0&\text{\lr{for }}0<\phi<\pi\\-V_0&\text{\lr{for }}\pi<\phi<2\pi\end{cases}}

\begin{itemize}
\tightlist
\item
  درون لوله (\ELim{r<b})

  \begin{itemize}
  \tightlist
  \item
    با توجه به اینکه ناحیه مورد نظر شامل \ELim{r=0} است، جملات شامل \ELim{r^{-n}} نمی‌توانند وجود داشته باشند.
  \item
    با توجه به اینکه \ELim{V(r,\phi)} یک تابع فرد از \ELim{\phi} است، بنابراین
  \end{itemize}
\end{itemize}

\ELdisp{V_n(r,\phi)\ELbr =A_nr^n\sin n\phi}

\begin{itemize}
\tightlist
\item
  عبارت فوق به تنهایی نمی تواند شرایط مرزی را برآورده کند. بنابراین
\end{itemize}

\ELdisp{\begin{aligned}V(r,\phi)&=\sum_{n=1}^{\infty}V_n(r,\phi)\\&=\sum_{n=1}^{\infty}A_nr^n\sin n\phi\end{aligned}}
\ELdisp{\sum_{n=1}^{\infty}A_nb^n\sin n\phi\ELbr =\begin{cases}V_0&\text{\lr{for }}0<\phi<\pi\\-V_0&\text{\lr{for }}\pi<\phi<2\pi\end{cases}}

\begin{itemize}
\tightlist
\item
  رابطه فوق بسط سری فوریه یک تابع متناوب و فرد است که در بازه \ELim{0<\phi<\pi} مقدار آن \ELim{V_0} و در بازه \ELim{\pi<\phi<2\pi} مقدار آن \ELim{-V_0} می باشد.
\end{itemize}

\ELdisp{A_n\ELbr =\begin{cases}\dfrac{4V_0}{n\pi b^n}&\text{\lr{if }}n\text{\lr{ is odd}}\\\noalign{\vskip 2mm}0&\text{\lr{if }}n\text{\lr{ is even}}\end{cases}\qquad\ELbr \qquad\ELbr  V(r,\phi)\ELbr =\frac{4V_0}{\pi}\sum_{n=\text{\lr{odd}}}^{\infty}\frac1n\left(\frac rb\right)^n\sin n\phi,\qquad\ELbr  r<b}

\begin{itemize}
\tightlist
\item
  بیرون لوله (\ELim{r>b})

  \begin{itemize}
  \tightlist
  \item
    با توجه به اینکه ناحیه مورد نظر شامل \ELim{r\to\infty} است، جملات شامل \ELim{r^n} نمی‌توانند وجود داشته باشند.
  \item
    \ELim{V(r,\phi)} یک تابع فرد از \ELim{\phi} است.
  \item
    یک عبارت به تنهایی نمی تواند شرایط مرزی را برآورده کند. بنابراین
  \end{itemize}
\end{itemize}

\ELdisp{\begin{aligned}V(r,\phi)&=\sum_{n=1}^{\infty}V_n(r,\phi)\\&=\sum_{n=1}^{\infty}B'_nr^{-n}\sin n\phi\end{aligned}\qquad\ELbr \begin{aligned}V(b,\phi)&=\sum_{n=1}^{\infty}B'_nb^{-n}\sin n\phi\\&=\begin{cases}V_0&\text{\lr{for }}0<\phi<\pi\\-V_0&\text{\lr{for }}\pi<\phi<2\pi\end{cases}\end{aligned}}
\ELdisp{B'_n\ELbr =\begin{cases}\dfrac{4V_0b^n}{n\pi}&\text{\lr{if }}n\text{\lr{ is odd}}\\\noalign{\vskip 2mm}0&\text{\lr{if }}n\text{\lr{ is even}}\end{cases}}

\begin{ELimportant}{}

\ELdisp{V(r,\phi)\ELbr =\frac{4V_0}{\pi}\sum_{n=\text{\lr{odd}}}^{\infty}\frac1n\left(\frac br\right)^n\sin n\phi,\qquad\ELbr  r>b}

\end{ELimportant}

\end{ELsolution}

\end{ELexample}
